\documentclass[11pt]{article}
\usepackage{amsmath,amsthm,amsfonts,amssymb,amscd}
\usepackage{multirow,booktabs}
\usepackage[table]{xcolor}
\usepackage{fullpage}
\usepackage{lastpage}
\usepackage{enumitem}
\usepackage{fancyhdr}
\usepackage{mathrsfs}
\usepackage{wrapfig}
\usepackage{setspace}
\usepackage{calc}
\usepackage{multicol}
\usepackage{cancel}
\usepackage[retainorgcmds]{IEEEtrantools}
\usepackage[margin=3cm]{geometry}
\usepackage{amsmath}
\newlength{\tabcont}
\setlength{\parindent}{0.0in}
\setlength{\parskip}{0.05in}
\usepackage{empheq}
\usepackage{framed}
% \usepackage{CJKutf8}
\usepackage[most]{tcolorbox}
\usepackage{xcolor}
\usepackage[]{titlesec}

% https://tex.stackexchange.com/questions/250052/custom-lstlisting-highlighting-for-symbols
\usepackage{listings}

\usepackage{tcolorbox}

\usepackage{quantikz}

\newcommand{\lstfont}[1]{\color{#1}\scriptsize\ttfamily}

\lstset{
    language=[ANSI]C++,
    showstringspaces=false,
    backgroundcolor=\color{black!90},
    basicstyle=\lstfont{white},
    identifierstyle=\lstfont{white},
    keywordstyle=\lstfont{magenta!40},
    frame=l,
    framesep=4.5mm,
    framexleftmargin=2.5mm,
    numbers=left,                    
    numbersep=5pt,
    numberstyle=\lstfont{white},
    stringstyle=\lstfont{cyan},
    commentstyle=\lstfont{yellow!30},
    emph={
        cudaMalloc, cudaFree,
        __global__, __shared__, __device__, __host__,
        __syncthreads,
    },
    emphstyle={\lstfont{green!60!white}},
    breaklines=true
}

\colorlet{shadecolor}{orange!15}
\parindent 0in
\parskip 12pt
\geometry{margin=1in, headsep=0.25in}
\theoremstyle{definition}
\newtheorem{defn}{Definition}
\newtheorem{reg}{Rule}
\newtheorem{exer}{Exercise}
\newtheorem{note}{Note}
% \renewcommand\thesection{Question \arabic{section}}
% \renewcommand\thesubsection{\arabic{section}.(\alph{subsection})}


\newcolumntype{P}[1]{>{\centering\arraybackslash}p{#1}}
\newcolumntype{M}[1]{>{\centering\arraybackslash}m{#1}}
\newcolumntype{R}[1]{>{\raggedleft\arraybackslash}m{#1}}

\begin{document}
%\setcounter{section}{0}
\title{Homework 5}

\thispagestyle{empty}

\begin{center}
    {\LARGE \bf UQSD Problem formulation}\\
    \bigskip
    \today\\
\end{center}

\section{Optimal POVM for linearly independent pure states}
% TODO \cite{}
Denote a finite set of quantum states $\{\ket{\psi_{1}}, \ket{\psi_{2}}, \dots, \ket{\psi_{k}}\}$,
where $k \in \mathbf{N}$, in finite size Hilbert space?, prepared in the prior probabilities
$\{p_{1}, p_{2}, \dots, p_{k}\}$,
and a set of measurement operators \dots
The reciprocal states defined in %\cite{}
can be calculated with \dots

By precalculating the reciprocal states, the problem can be reduced to find a set of $k$ parameters that
while keeping the operators positive semidefinite, maximizes the total joint success probability.
These parameters actually correspond to the success probability of each outcome.
The semidefinite programming problem can be formulated as follows:
\begin{align*}
    7788
\end{align*}

\section{Fixed rate of inconclusive outcome for mixed-state discrimination}
%\cite{origin} \cite{application} \cite{follow-ups}
A quantum state $\Psi$ can be denoted by a state vector $\ket{\psi}$ or a density matrix $\rho$.
The prior probability for each input state $\Psi_{i}$ (expressed by $\ket{\psi_{i}}$ or $\rho_{i}$) is
denoted by $p_{i}$, where $i$ ranges from 1 to $k$.
The measurement operators are defined by $\Pi_{i}$, where $i$ ranges from 0 to $k$.
Here $\Pi_{0}$ is designed to match the inconclusive results,
and $\Pi_{1}$ to $\Pi_{k}$ designed to measure the $1$-st to $k$-th state respectively.
$P_{D}$ denotes the total probability of correct discrimination,
$P_{err}$ denotes the probability of incorrect/erroneous discrimination,
and $P_{inc}$ denotes the probability of inconclusive outcomes.
$\beta$ is a constant real number that is defined to be the target inconclusive probability.
In practice, we use an inequality in place of an equality,
and $\beta$ is the minimum rate of inconclusive outcomes.

\begin{align}
    P_{D} + P_{err} + P_{inc} = 1
\end{align}

\begin{align}
     & \underset{\{\Pi\}}{\text{maximize}} &  & P_{D}=\sum_{i=1}^{k} p_{i} \operatorname{Tr}(\rho_{i} \Pi_{i})              \\
     & \text{subject to}                   &  & \Pi_{i} \succeq 0, \quad i = 0, \ldots, k                                   \\
     &                                     &  & \sum_{i=0}^{k} \Pi_i = I                                                    \\
     &                                     &  & P_{inc}=\sum_{i=1}^{k} p_{i} \operatorname{Tr}(\rho_{i} \Pi_{0}) \geq \beta
\end{align}

\section{$\epsilon$-UQSD}
In this work, we include a different set of constraints in the semidefinite programming problem
such that a desired alternative POVM can be obtained.
% Assume a set of $k$ input quantum states are expressed in density matrices,
% denoted by $\{ \rho_{1}, \rho_{2}, \dots, \rho_{k}\}$.
The notations are almost identical to FRIO, but our constraints further consider the error part of the conclusive results.
It also resembles the max confidence discrimination scheme (TBA).
The symbol $O_{i}$ denotes the event that the $i$-th outcome is measured.
The density matrices could be treated without assuming that the states are related to some pure states.
Later, we may assume the mixed states are related to some pure states or can be expressed by ensembles of pure states.
The semidefinite programming formulation is as follows:

\begin{align}
     & \underset{\{\Pi\}}{\text{maximize}} &  & P_{D}=\sum_{i=1}^{k} p_{i} \operatorname{Tr}(\rho_{i} \Pi_{i})                                                                                                      \\
     & \text{subject to}                   &  & \Pi_{i} \succeq 0, \quad i = 0, \ldots, k                                                                                                                           \\
     &                                     &  & \sum_{i=0}^{k} \Pi_i = I                                                                                                                                            \\
     &                                     &  & P(O_i \mid \rho_i) = \frac{p_i \, \operatorname{Tr}(\rho_i \Pi_i)}{\sum_{j=1}^{L} p_i \, \operatorname{Tr}(\rho_i \Pi_j)} \geq 1 - \alpha_i, \quad i = 1, \ldots, k \\
     &                                     &  & P(\rho_i \mid O_i) = \frac{p_i \, \operatorname{Tr}(\rho_i \Pi_i)}{\sum_{j=1}^{L} p_j \, \operatorname{Tr}(\rho_j \Pi_i)} \geq 1 - \beta_i, \quad i = 1, \ldots, k
\end{align}
% &  & \gamma_{i} \ge \sum_{j=1, j \neq i}^{m} \rho_{i} \operatorname{Tr}(\rho_{i} \Pi_{j} ), 1≤𝑖≤𝑚, 𝛾∈𝑅^𝑚            \\
% &  & P_{err} = \sum_{i=1}^{m} \sum_{j=1, j \neq i}^{m} \rho_{i} \operatorname{Tr}(\rho_{i} \Pi_{j} ) \label{eq:p_err} \\




\section{Discrimination for multiple non-orthogonal quantum states}
\section{Problem formulation}

% TODO Rewrite old stuff from last year
Suppose we have three 3-qubit states and we want to find a set of POVMs to
\begin{align*}
     & E^{\perp}_{1} \equiv |\psi^{\perp}_{[1]} \rangle \langle \psi^{\perp}_{[1]} |                                   \\
     & Given: \; |\psi_{i} \rangle \; \forall i \in \{ 1 .. k \}                                                       \\
     & Variables: \; |\psi^{\perp}_{[i]} \rangle \; \forall i \in \{ 0 .. k \}                                         \\
     & Minimize \; \Sigma_{j=1}^{k} \: p_{j} \langle \psi_{j}|  E^{\perp}_{0} |\psi_{j} \rangle                        \\
     & (1) \; \langle \psi_{j}|  E^{\perp}_{i} |\psi_{j} \rangle = 0 \; \forall i, j \in \{ 1 .. k \} \land (i \neq j) \\
     & (2) \; \langle \psi^{\perp}_{[i]} |\psi^{\perp}_{[j]} \rangle  == \delta_{ij} \; \forall i, j \in \{ 1 .. k \}  \\
     & (3) \; \Sigma_{i=0}^{k} E^{\perp}_{i} == I                                                                      \\
\end{align*}

\end{document}
